\documentclass[a4paper, 12pt]{article}

\usepackage[utf8]{inputenc}
\usepackage[T2A]{fontenc}
\usepackage[russian]{babel}
\usepackage{amsmath, amssymb, amsthm, amsfonts}
\usepackage{geometry}
\usepackage{algorithm}
\usepackage{algpseudocode}
\usepackage{graphicx}

\geometry{top=2cm, bottom=2cm, left=2.5cm, right=1.5cm}

\theoremstyle{definition}
\newtheorem{definition}{Определение}
\theoremstyle{plain}
\newtheorem{theorem}{Теорема}

\floatname{algorithm}{Алгоритм}
\renewcommand{\algorithmicrequire}{\textbf{Вход:}}
\renewcommand{\algorithmicensure}{\textbf{Выход:}}

\title{Билет №86 \\ \small Классическое определение систем. Импульсная характеристика системы. Многомерные системы, инвариантные к сдвигу. (СпБПУ)}
\author{Сериков Владислав Александрович}
\date{16 мая 2026}

\begin{document}

\maketitle
\tableofcontents
\newpage

\section{Классическое определение систем}

В рамках классической теории сигналов и систем под \textbf{системой} понимается математический оператор, который преобразует один сигнал (входной) в другой сигнал (выходной). Сигналы могут быть непрерывными (зависящими от непрерывного времени $t$) или дискретными (зависящими от дискретного отсчета $n$).

\begin{definition}
\textbf{Система} — это правило или оператор $\mathcal{H}$, устанавливающее однозначное соответствие между входным сигналом $x(t)$ (или $x[n]$) и выходным сигналом $y(t)$ (или $y[n]$):
$$y(t) = \mathcal{H}\{x(t)\}$$
Выходной сигнал $y(t)$ часто называют \textit{реакцией} или \textit{откликом} системы.
\end{definition}

Для того чтобы систему можно было эффективно анализировать математическими методами, выделяют ряд базовых свойств:

1. \textbf{Линейность.} Система называется линейной, если для нее выполняется принцип суперпозиции. Для любых входных сигналов $x_1(t)$ и $x_2(t)$ и любых комплексных констант $a, b$:
$$ \mathcal{H}\{a x_1(t) + b x_2(t)\} = a \mathcal{H}\{x_1(t)\} + b \mathcal{H}\{x_2(t)\} $$

2. \textbf{Инвариантность к сдвигу (стационарность).} Система инвариантна ко времени (сдвигу), если временной сдвиг входного сигнала приводит к такому же временному сдвигу выходного сигнала без изменения его формы. Если $y(t) = \mathcal{H}\{x(t)\}$, то для любого $t_0$:
$$ \mathcal{H}\{x(t - t_0)\} = y(t - t_0) $$

Системы, обладающие одновременно свойствами 1 и 2, называются \textbf{линейными стационарными системами} (ЛСС, англ. LTI -- Linear Time-Invariant, или LSI -- Linear Shift-Invariant).

3. \textbf{Причинность (казуальность).} Система называется причинной, если ее реакция в момент времени $t_0$ зависит только от значений входного сигнала при $t \le t_0$.

4. \textbf{Устойчивость.} Система называется устойчивой в смысле ОВ-ОР (ограниченный вход — ограниченная реакция, BIBO), если любой ограниченный входной сигнал $|x(t)| \le M_x < \infty$ вызывает ограниченный выходной сигнал $|y(t)| \le M_y < \infty$.

\section{Импульсная характеристика системы}

Особое значение в теории ЛСС имеет реакция системы на единичный импульс. Для непрерывных систем таким импульсом выступает дельта-функция Дирака $\delta(t)$, а для дискретных — дельта-функция Кронекера $\delta[n]$.

\begin{definition}
\textbf{Импульсная характеристика (ИХ)} $h(t)$ (или $h[n]$) — это реакция линейной стационарной системы, находящейся в нулевых начальных условиях, на входное воздействие в виде дельта-функции:
$$ h(t) = \mathcal{H}\{\delta(t)\}, \quad \text{или} \quad h[n] = \mathcal{H}\{\delta[n]\} $$
\end{definition}

Знание импульсной характеристики ЛСС позволяет полностью описать поведение системы. Это фундаментальное свойство выражается следующей теоремой.

\begin{theorem}[Интеграл свертки для непрерывных ЛСС]
Пусть $\mathcal{H}$ — линейная система, инвариантная к сдвигу, с импульсной характеристикой $h(t)$. Тогда для произвольного входного сигнала $x(t)$ реакция системы $y(t)$ полностью определяется интегралом свертки входного сигнала и импульсной характеристики:
$$ y(t) = \int_{-\infty}^{\infty} x(\tau) h(t - \tau) d\tau = x(t) * h(t) $$
\end{theorem}

\begin{proof}
Доказательство базируется на фильтрующем свойстве дельта-функции Дирака, согласно которому любой сигнал $x(t)$ можно представить как бесконечную сумму (интеграл) сдвинутых и взвешенных дельта-функций:
$$ x(t) = \int_{-\infty}^{\infty} x(\tau) \delta(t - \tau) d\tau $$

Применим к обеим частям оператор системы $\mathcal{H}$:
$$ y(t) = \mathcal{H}\{x(t)\} = \mathcal{H}\left\{ \int_{-\infty}^{\infty} x(\tau) \delta(t - \tau) d\tau \right\} $$

Так как система линейная, оператор $\mathcal{H}$ можно внести под знак интеграла (поскольку интеграл является пределом линейной комбинации, а $\mathcal{H}$ линейна). Кроме того, $x(\tau)$ относительно переменной $t$ является константой:
$$ y(t) = \int_{-\infty}^{\infty} x(\tau) \mathcal{H}\{\delta(t - \tau)\} d\tau $$

Поскольку система инвариантна к сдвигу, реакция на сдвинутый импульс $\delta(t - \tau)$ есть сдвинутая импульсная характеристика $h(t - \tau)$:
$$ \mathcal{H}\{\delta(t - \tau)\} = h(t - \tau) $$

Подставляя это выражение в интеграл, получаем:
$$ y(t) = \int_{-\infty}^{\infty} x(\tau) h(t - \tau) d\tau $$
Теорема доказана.
\end{proof}

Для дискретных систем аналогично используется \textbf{сумма свертки}:
$$ y[n] = \sum_{k=-\infty}^{\infty} x[k] h[n - k] $$

\subsection{Алгоритм дискретной свертки}

При программной реализации систем с конечной импульсной характеристикой (КИХ-фильтров) используется алгоритм вычисления свертки для массивов конечной длины.

\begin{algorithm}[H]
\caption{Вычисление дискретной свертки двух конечных сигналов}
\begin{algorithmic}[1]
\Require Массив сигнала $x$ длины $N$, массив ИХ $h$ длины $M$.
\Ensure Массив выходного сигнала $y$ длины $L = N+M-1$.
\State Выделить память под массив $y$ длины $L$ и заполнить его нулями.
\For{$i \gets 0$ \textbf{to} $N-1$}
    \For{$j \gets 0$ \textbf{to} $M-1$}
        \State $y[i+j] \gets y[i+j] + x[i] \cdot h[j]$
    \EndFor
\EndFor
\State \Return $y$
\end{algorithmic}
\end{algorithm}

\textbf{Иллюстрация алгоритма на минимальном примере:}
Пусть входной сигнал $x = [1, 2]$ (длина $N=2$), а ИХ системы $h = [3, -1]$ (длина $M=2$).
Размер выхода $L = 2 + 2 - 1 = 3$. Инициализация: $y = [0, 0, 0]$.
\begin{itemize}
    \item \textbf{Шаг 1 ($i=0, x[0]=1$):}
    \begin{itemize}
        \item $j=0, h[0]=3$: $y[0+0] = y[0] + 1 \cdot 3 \implies y = [3, 0, 0]$
        \item $j=1, h[1]=-1$: $y[0+1] = y[1] + 1 \cdot (-1) \implies y = [3, -1, 0]$
    \end{itemize}
    \item \textbf{Шаг 2 ($i=1, x[1]=2$):}
    \begin{itemize}
        \item $j=0, h[0]=3$: $y[1+0] = y[1] + 2 \cdot 3 \implies y = [3, -1+6, 0] = [3, 5, 0]$
        \item $j=1, h[1]=-1$: $y[1+1] = y[2] + 2 \cdot (-1) \implies y = [3, 5, -2]$
    \end{itemize}
\end{itemize}
\textbf{Ответ:} $y = [3, 5, -2]$.

\section{Многомерные системы, инвариантные к сдвигу}

Понятие систем легко обобщается на случай сигналов, зависящих от нескольких независимых переменных. Наиболее частый пример в инженерии — двумерные (2D) сигналы, такие как изображения, где независимыми переменными являются пространственные координаты.

Многомерный сигнал $x(n_1, n_2, \dots, n_D)$ подвергается преобразованию многомерным оператором $\mathcal{H}$.

\begin{definition}
\textbf{Двумерная линейная система, инвариантная к пространственному сдвигу (2D ЛСИС)}, удовлетворяет двум условиям:
\begin{enumerate}
    \item $\mathcal{H}\{a x_1[n_1, n_2] + b x_2[n_1, n_2]\} = a\mathcal{H}\{x_1\} + b\mathcal{H}\{x_2\}$
    \item Если $y[n_1, n_2] = \mathcal{H}\{x[n_1, n_2]\}$, то для любых целых сдвигов $m_1, m_2$:
    $$ \mathcal{H}\{x[n_1 - m_1, n_2 - m_2]\} = y[n_1 - m_1, n_2 - m_2] $$
\end{enumerate}
\end{definition}

Двумерный единичный импульс (дельта-функция) определяется как:
$$
\delta[n_1, n_2] =
\begin{cases}
1, & \text{если } n_1 = 0 \text{ и } n_2 = 0 \\
0, & \text{в противном случае}
\end{cases}
$$

Импульсная характеристика 2D-системы представляет собой двумерный массив (или матрицу) значений: $h[n_1, n_2] = \mathcal{H}\{\delta[n_1, n_2]\}$. В обработке изображений матрица $h[n_1, n_2]$ часто называется \textbf{ядром свертки} (convolution kernel) или \textbf{маской фильтра}.

\begin{theorem}[Двумерная дискретная свертка]
Выходной сигнал двумерной ЛСИС выражается через двумерную сумму свертки входного сигнала $x[n_1, n_2]$ и импульсной характеристики $h[n_1, n_2]$:
$$ y[n_1, n_2] = x[n_1, n_2] * h[n_1, n_2] = \sum_{k_1=-\infty}^{\infty} \sum_{k_2=-\infty}^{\infty} x[k_1, k_2] h[n_1 - k_1, n_2 - k_2] $$
\end{theorem}
\begin{proof}
Свойства линейности и инвариантности к сдвигу в многомерном пространстве применяются аналогично одномерному случаю. Любой 2D сигнал можно разложить по базису сдвинутых импульсов:
$$ x[n_1, n_2] = \sum_{k_1} \sum_{k_2} x[k_1, k_2] \delta[n_1 - k_1, n_2 - k_2] $$
В силу линейности $\mathcal{H}$ вносится под двойную сумму, воздействуя только на $\delta$. В силу пространственной инвариантности $\mathcal{H}\{\delta[n_1 - k_1, n_2 - k_2]\} = h[n_1 - k_1, n_2 - k_2]$, откуда напрямую следует формула двумерной свертки.
\end{proof}

\textbf{Минимальный пример 2D ЛСИС (Размытие изображения / Фильтр усреднения):}
Пусть ядро фильтра $h$ задано матрицей $3 \times 3$, где каждый элемент равен $1/9$. Эта система заменяет значение каждого пикселя $x[n_1, n_2]$ на среднее арифметическое его соседей.
Если на вход подается изображение с резким перепадом яркости, импульсная характеристика "размажет" этот перепад, осуществив пространственную низкочастотную фильтрацию (эффект "блюра").

\end{document}
